\begin{resumo}[\protect\bfseries Resumo]
Este memorial descritivo apresenta a trajetória de Arthur Fenilli de Souza nas quatro etapas da Agência Experimental de Engenharia de Software da Pontifícia Universidade Católica do Rio Grande do Sul, entre 2023 e 2026. O documento reúne o contexto dos projetos, as tecnologias utilizadas, as atividades desenvolvidas em cada sprint e uma reflexão sobre os aprendizados técnicos e pessoais. Na AGES I, o projeto Plan Line proporcionou o primeiro contato com o desenvolvimento de uma aplicação em equipe e com demandas de clientes. Na AGES II, o projeto de Certificação por meio de Blockchain aprofundou a experiência com modelagem de dados, backend e operações assíncronas. Na AGES III, o VínculoHub Portal ampliou a participação em decisões técnicas, organização de tarefas e acompanhamento de colegas. Na AGES IV, o ReadRace reuniu responsabilidades de gestão, desenvolvimento, revisão de código, testes e apresentação das entregas. A trajetória também contempla a evolução do uso de agentes de inteligência artificial e a relação entre os conhecimentos das disciplinas e a prática dos projetos. A reflexão destaca a importância do planejamento, da comunicação e da colaboração e sua influência nas primeiras experiências profissionais e na direção de carreira do autor.

\vspace{\onelineskip}
\noindent\textbf{Palavras-chave:} Engenharia de Software. AGES. Desenvolvimento de software. Trabalho em equipe. Formação profissional.
\end{resumo}
